\chapter{Preface}

This book grew out of a paper presented at the workshop on data security and
privacy that Banco de Portugal and the Irving Fisher Committee held in Lisbon in
September 2026. The paper argued that a disclosure risk figure only becomes useful
to an institution once it is attached to a declared attack, an honest statement of
its uncertainty and a written decision rule. In discussion, a participant asked
what happens to the risk when an institution gives several data sets to the same
researcher. The question went to the centre of the argument, and a conference
paper could only answer part of it. Chapter~\ref{ch:programmes} answers it in
full, and the rest of the book builds the framework that answer needs.

The book has two audiences, and it is arranged so that each can read it without
the other's material getting in the way. The first is the statistician, data
protection officer or release committee member at a central bank, statistical
office or health registry who has to decide whether a file goes out. For that
reader every chapter opens with the problem in institutional terms, states its
results in words before it states them in symbols, and closes its main text with a
short block headed \textit{In practice} that says what to write down and what to
check. The second audience is the researcher in statistics, computer science or
economics who wants to know exactly what is claimed and under which assumptions.
For that reader every result carries a numbered statement and a full proof, and
the proofs sit in a final section of each chapter so that the practitioner can
pass over them without losing the thread.

Every number in the book comes from one research pipeline, the Python
package \texttt{privacy\_risk} in the repository that accompanies the book. Where
a chapter reports a figure, it names the module and function that produced it, so
that a reader can rerun the calculation, change an assumption and see what moves.
Most of the evidence is synthetic. That is a limitation and the book says so
in the chapters it affects: the credit register in Part~V shares the structure of a
supervisory collection, but its parameters are chosen, not estimated, and the
floors it implies are conditional on that design.

The book also reports what did not work. The chapter on model-proposed attacks
includes a pilot in which augmenting the training data did not improve the model,
and the audit chapter shows how often an analyst who retries a failed audit
certifies an unsafe release. These results belong in the record because a method
for governing risk has to be judged on its failures as well as its successes.

This draft holds the preface, the first two chapters and Chapter~9. The other
chapters appear as short plans so that the reader can see where the argument goes.
Comments on the draft are welcome.
